\documentclass[paper=a3paper,landscape]{standalone}
\usepackage{fouriernc}
\usepackage{trfsigns}
\usepackage{pgfplots}
\usetikzlibrary{calc}
\usepackage{comment}
\usepackage{url}
\usepackage{amssymb}
\usepackage{xfrac}
\usepackage{sig_sys_macros}
\newcommand{\sinc}{\mathrm{sinc}}
\newcommand{\rect}{\mathrm{rect}}
\newcommand{\psinc}{\mathrm{psinc}}
% uncomment for quadratic paper
%\renewcommand{\paperwidth}{\paperheight}
\begin{document}
\begin{tikzpicture}[
declare function = {
	sinc(\t)=sin(deg(\t))/\t*(\t!=0)+1*(\t==0);
	rect(\t)=and(\t>=-1/2,\t<=1/2);
	square(\t)= and(mod(\t,2)>=0,mod(\t,2)<1);
	RECT(\t,\n) = 1*and(\t>=0,\t<\n);
	diric(\w,\n) =( mod(abs(\w),2*pi)==0)*\n +( mod(abs(\w),2*pi)!=0)*sin(deg(\n/2*abs(\w)))/sin(deg(abs(\w)/2));
	diric2(\u,\m,\n)=( mod(abs(\u),\n)==0)*\m +(mod(abs(\u),\n)!=0)*sin(deg(\m/\n*pi*\u))/sin(deg(1/\n*pi*\u));
},
]
% https://tex.stackexchange.com/questions/184111/how-do-i-make-label-node-to-appear-in-front-of-markers-in-pgfplots
% ticks labels in front of addplot:
\pgfplotsset{
  layers/mylayers/.define layer set=
    {axis background,axis grid,axis ticks,axis lines,%
      main,axis tick labels,axis descriptions,axis foreground}
    {/pgfplots/layers/standard}% re-uses the style definitions of layers/standard
}
\pgfplotsset{set layers=mylayers} % using mylayers
%
% period
\newcommand{\T}{5}
% impulse width
\newcommand{\Th}{2}
% axis scaling below is hard coded!
% label font size
\newcommand{\FONTSIZE}{\scriptsize}
% time domain color
\newcommand{\TCOLOR}{C0}
% frequency domain color
\newcommand{\FCOLOR}{C1}
% line width of continuous functions
\newcommand{\LINEWIDTH}{ultra thick}
% marker size of discret functions
\newcommand{\MARKSIZE}{1.6pt}
\pgfmathsetmacro{\W}{pi}
\pgfmathsetmacro{\PAPERWIDTHSQUARED}{\paperwidth/32/32*\paperwidth}
\pgfmathsetmacro{\PAPERHEIGHTSQUARED}{\paperheight/32/32*\paperheight}
\pgfmathsetmacro{\ANGLE}{atan(\paperheight/\paperwidth)}
%
%
%
\draw[C7, ultra thick] (0,0) rectangle (\paperwidth,\paperheight);
%
%
%
% red cross
%
\draw[C7,ultra thick] (3/8*\paperwidth,1/2*\paperheight) -- (5/8*\paperwidth,1/2*\paperheight);
\draw[C7,ultra thick] (1/2*\paperwidth,3/8*\paperheight) -- (1/2*\paperwidth,5/8*\paperheight);
%
%
%
% continuation of cross lines at the paper border
%
\draw[C7,ultra thick] (0,1/2*\paperheight) -- (1/8*\paperwidth,1/2*\paperheight);
\draw[C7,ultra thick] (7/8*\paperwidth,1/2*\paperheight) -- (\paperwidth,1/2*\paperheight);
\draw[C7,ultra thick] (1/2*\paperwidth,1/64*\paperheight) -- (1/2*\paperwidth,1/8*\paperheight);
\draw[C7,ultra thick] (1/2*\paperwidth,7/8*\paperheight) -- (1/2*\paperwidth,\paperheight);
%
%
%
% transformation abbreviations in corners
%
\node[anchor=north west] at (0,\paperheight) (FT) {\Huge {FT} \small Fourier Transform};
\node[anchor=north east] at (\paperwidth,\paperheight) (FS) {\small Fourier Series \Huge {FS}};
\node[anchor=south west] at (0,0) (DTFT) {\Huge {DTFT} \small Discrete-Time Fourier Transform};
\node[anchor=south east] at (\paperwidth,0) (DFT) {\small Discrete Fourier Transform \Huge {DFT}};
%
\node[anchor=north west] at (11/64*\paperwidth,17/20*\paperheight) (wcFT) {\FONTSIZE $\omega_c = \frac{2\pi}{T_c}$};
\node[anchor=north west] at (25/32*\paperwidth,17/20*\paperheight) (woFS) {\FONTSIZE $\omega_0 = \frac{2\pi}{T_0}$};
\node[anchor=north west] at (25/32*\paperwidth,5/20*\paperheight) (NFFT) {\FONTSIZE $N=10$};
\node[anchor=north west] at (24.5/32*\paperwidth,4.5/20*\paperheight) (NFFTnote) {\FONTSIZE exact DFT pair};
\node[anchor=north west] at (11/64*\paperwidth,5/20*\paperheight) (WcDTFT) {\FONTSIZE $\Omega_c = \frac{2\pi}{K_c}$};
\node[anchor=north west] at (11/64*\paperwidth,4.5/20*\paperheight) (WcDTFTfor) {\FONTSIZE here for $K_c=2$};
%

% FT
\node[anchor=north west] at (1/64*\paperwidth,19/20*\paperheight) (FTequation) {$
x(t)=\frac{1}{2 \pi} \int\limits_{\omega=-\infty}^{+\infty} X(\omega) \e^{+\im \omega t} \fsd \omega
\qquad
X(\omega)=\int\limits_{t=-\infty}^{+\infty} x(t) \e^{-\im \omega t}\fsd t
\qquad
t\in\mathbb{R}, \omega\in\mathbb{R}
$};
%DTFT
\node[anchor=north west] at (1/64*\paperwidth,2/20*\paperheight) (DTFTequation) {$
x[k]=\frac{1}{2 \pi} \int\limits_{\Omega=-\pi}^{+\pi} X(\Omega) \e^{+\im \Omega k}
\qquad
X(\Omega)=\sum\limits_{k=-\infty}^{+\infty} x[k] \e^{-\im \Omega k}
\qquad
k\in\mathbb{Z}, \Omega\in\mathbb{R}
$};
%DFT
\node[anchor=north west] at (40/64*\paperwidth,2/20*\paperheight) (DFTequation) {$
x[k]=\frac{1}{N} \sum\limits_{\mu=0}^{N-1} X[\mu] \e^{+\im \frac{2 \pi}{N} \mu k}
\qquad
X[\mu]=\sum\limits_{k=0}^{N-1} x[k] \e^{-\im \frac{2 \pi}{N} \mu k}
\qquad
k\in\mathbb{Z}, \mu\in\mathbb{Z}
$};
%FS
\node[anchor=north west] at (38.5/64*\paperwidth,19.2/20*\paperheight) (FSequation) {$
x(t)=\frac{1}{T_0} \sum\limits_{\mu=-\infty}^{+\infty} \tilde{X}[\mu] \e^{+\im \frac{2 \pi}{T_0} \mu t}
\qquad
\tilde{X}[\mu]=\int\limits_{t=-\frac{T_0}{2}}^{+\frac{T_0}{2}} x(t) \e^{-\im \frac{2 \pi}{T_0} \mu t} \fsd t
\qquad
t\in\mathbb{R}, \mu\in\mathbb{Z}
$};

\node[anchor=north west] at (40/128*\paperwidth,2/128*\paperheight) (author) {\tiny Robert Hauser, Frank Schultz, University of Rostock, \url{https://github.com/spatialaudio/signals-and-systems-cheatsheet}, CC BY 4.0, 2026-05-19};
% FT Plots
%
%
% sinc time domain
%
\begin{axis}[
	axis lines = middle,
	axis line style = {-latex},
	xlabel = {$t$},
	x label style = {anchor=west},
	ylabel = {$x(t)$},
	y label style={anchor=south},
	width = 18/80*\paperwidth,
	height = 18/80*\paperheight,
	at = {(1/10*\paperwidth,14/20*\paperheight)},
	xtick = {-7,-6,-5,-4,-3,-2,-1,0,1,2,3,4,5,6,7},
	xticklabels = {,,$-5T_c$,,,,$-T_c$,,$T_c$,,,,$5T_c$,,},
	ytick = {0,0.5,1},
	yticklabels = {},
	ticklabel style={font=\FONTSIZE},
	grid = major,
	grid style={C0!50},
	ymin = -0.25,
	ymax = 1.125,
	xmin = -8.5,
	xmax = 8.5,
	anchor = south,
	]
	\addplot[domain=-7.5:7.5,
	samples = 257,
	color=\TCOLOR,
	\LINEWIDTH,]
	{sinc(x*\W)};
	\node[anchor=center] at (axis cs:-0.75,1) {\FONTSIZE $A$};
\end{axis}
%
% sinc frequency domain
%
\begin{axis}[
	axis lines = center,
	axis line style = {-latex},
	xlabel = {$\omega$},
	x label style = {anchor=west},
	ylabel = {$X\left(\omega\right)$},
	y label style={anchor=south},
	width = 18/80*\paperwidth,
	height = 18/80*\paperheight,
	at = {(3/10*\paperwidth,14/20*\paperheight)},
	xtick = {-pi,0,pi},
	xticklabels = {$\frac{-\omega_c}{2}$,,$\frac{\omega_c}{2}$},
	ytick = {0,0.5,1},
	yticklabels = {},
	ticklabel style={font=\FONTSIZE},
	grid=major,
	grid style={C1!40},
	ymin = -0.25,
	ymax = 1.125,
	xmin = -17/15*6*pi,
	xmax = 17/15*6*pi,
	anchor=south,
	]
	\addplot[domain=-6*pi:6*pi,
	samples = 257,
	color=\FCOLOR,
	\LINEWIDTH,]
	{pi/(\W)*rect(x/2*1/(\W))};
	\node[anchor=center] at (axis cs:-5*2*pi/\T,1) {\FONTSIZE $A T_c$};
\end{axis}
%
%
%
% FS Plots
%
%
% sinc time domain
%
\begin{axis}[
	axis lines = middle,
	axis line style = {-latex},
	xlabel = {$t$},
	x label style = {anchor=west},
	ylabel = {$x(t)$},
	y label style={anchor=south},
	width = 18/80*\paperwidth,
	height = 18/80*\paperheight,
	at = {(7/10*\paperwidth,14/20*\paperheight)},
	xtick = {-7.5,-5,...,7.5},
	xticklabels = {,$-T_0$,,,,$T_0$,},
	ytick = {0,0.5,1},
	yticklabels = {},
	ticklabel style={font=\FONTSIZE},
	grid = major,
	grid style={C0!50},
	anchor = south,
	ymin = -0.25,
	ymax = 1.125,
	xmin = -8.5,
	xmax = 8.5,
	]
	\addplot[domain=-7.5:7.5,
	samples = 257,
	color=\TCOLOR,
	\LINEWIDTH,]
	{diric(x/5*2*pi,5)/5};
	\node[anchor=center] at (axis cs:-0.75,1) {\FONTSIZE $A$};
\end{axis}
%
% sinc frequency domain
%
\begin{axis}[
	axis lines = center,
	axis line style = {-latex},
	xlabel = {$\mu$},
	x label style = {anchor=west},
	ylabel = {$\tilde{X}\left[\mu\right]$},
	y label style={anchor=south},
	width = 18/80*\paperwidth,
	height = 18/80*\paperheight,
	at = {(9/10*\paperwidth,14/20*\paperheight)},
	xtick = {-15,-10,-5,0,5,10,15},
	xticklabels = {,,$-5$,,$5$,$10$,},
	ytick = {0,0.5,1},
	yticklabels = {},
	ticklabel style={font=\FONTSIZE},
	grid = major,
	grid style={C1!40},
	ymin = -0.25,
	ymax = 1.125,
	xmin = -17,
	xmax = 17,
	anchor=south,
	]
	\addplot[domain=-15:15,
	samples = 31,
	color=\FCOLOR,
	ycomb,
	mark=*,
	mark options = {scale=0.75},]
	{rect(x/4)};
	\node[anchor=center] at (axis cs:-3,1) {\FONTSIZE $\frac{A}{5}$};
\end{axis}
%
%
%
%
% DTFT Plots
%
%
% sinc time domain
%
\begin{axis}[
	axis lines = middle,
	axis line style = {-latex},
	xlabel = {$k$},
	x label style = {anchor=west},
	ylabel = {$x[k]$},
	y label style={anchor=south},
	width = 18/80*\paperwidth,
	height = 18/80*\paperheight,
	at = {(1/10*\paperwidth,6/20*\paperheight)},
	xtick = {-15,-10,-5,0,2,5,10,15},
	xticklabels = {,,$-5$,,$K_c$,$5$,$10$,},
	ytick = {0,0.5,1},
	yticklabels = {},
	ticklabel style={font=\FONTSIZE},
	grid = major,
	grid style={C0!50},
	ymin = -0.25,
	ymax = 1.125,
	xmin = -17,
	xmax = 17,
	anchor = north,
	]
	\addplot[domain=-15:15,
	samples = 31,
	color=\TCOLOR,
	ycomb,
	mark=*,
	mark size = \MARKSIZE,]
	{sinc(x/2*pi)};
	\node[anchor=center] at (axis cs:-1,1) {\FONTSIZE $A$};
\end{axis}
%
% sinc frequency domain
%
\begin{axis}[
	axis lines = center,
	axis line style = {-latex},
	xlabel = {$\Omega$},
	x label style = {anchor=west},
	ylabel = {$X\left(\Omega\right)$},
	y label style={anchor=south},
	width = 18/80*\paperwidth,
	height = 18/80*\paperheight,
	at = {(3/10*\paperwidth,6/20*\paperheight)},
	xtick = {-3*pi,-2*pi,-1*pi,-pi/2,0,pi/2,1*pi,2*pi,3*pi},
	xticklabels = {,$-2\pi$,,$\frac{-\Omega_c}{2}$,,$\frac{\Omega_c}{2}$,,$2\pi$,},
	ytick = {0,1,2},
	yticklabels = {},
	ticklabel style={font=\FONTSIZE},
	grid = major,
	grid style={C1!40},
	ymin = -0.5,
	ymax = 2.25,
	xmin = -17/15*3*pi,
	xmax = 17/15*3*pi,
	anchor=north,
	]
	\addplot[domain=-3*pi:3*pi,
	samples = 257,
	color=\FCOLOR,
	\LINEWIDTH,]
	{2*rect(x/pi)+2*rect((x-2*pi)/pi)+2*rect((x+2*pi)/pi)};
	\node[anchor=center] at (axis cs:-pi,2) {\FONTSIZE $A K_c$};
\end{axis}
%
%
%
% DFT Plots
%
%
% sinc time domain
%
\begin{axis}[
	axis lines = middle,
	axis line style = {-latex},
	xlabel = {$k$},
	x label style = {anchor=west},
	ylabel = {$x[k]$},
	y label style={anchor=south},
	width = 18/80*\paperwidth,
	height = 18/80*\paperheight,
	at = {(7/10*\paperwidth,6/20*\paperheight)},
	xtick = {-15,-10,-5,0,5,10,15},
	xticklabels = {,,$-5$,,$5$,$10$,},
	ytick = {0,0.5,1},
	yticklabels = {},
	ticklabel style={font=\FONTSIZE},
	grid = major,
	grid style={C0!50},
	ymin = -0.25,
	ymax = 1.125,
	xmin = -17,
	xmax = 17,
	anchor = north,
	]
	\addplot[domain=-15:15,
	samples = 31,
	color=\TCOLOR,
	ycomb,
	mark=*,
	mark size = \MARKSIZE,]
	{diric(x*2*pi/10,5)/5};
	\node[anchor=center] at (axis cs:-1,1) {\FONTSIZE $A$};
\end{axis}
%
% sinc frequency domain
%
\begin{axis}[
	axis lines = center,
	axis line style = {-latex},
	xlabel = {$\mu$},
	x label style = {anchor=west},
	ylabel = {$X\left[\mu\right]$},
	y label style={anchor=south},
	width = 18/80*\paperwidth,
	height = 18/80*\paperheight,
	at = {(9/10*\paperwidth,6/20*\paperheight)},
	xtick = {-15,-10,-5,0,5,10,15},
	xticklabels = {,,$-5$,,$5$,$10$,},
	ytick = {0,1,2},
	yticklabels = {},
	ticklabel style={font=\scriptsize},
	grid = major,
	grid style={C1!40},
	ymin = -0.5,
	ymax = 2.25,
	xmin = -17,
	xmax = 17,
	anchor=north,
	]
	\addplot[domain=-15:15,
	samples = 31,
	color=\FCOLOR,
	ycomb,
	mark = *,
	mark size = \MARKSIZE,]
	{RECT(x+2,5)*2+RECT(x+2+\T*2,5)*2+RECT(x+2-\T*2,5)*2};
	\node[anchor=center] at (axis cs:-4,2) {\FONTSIZE $2 A$};
\end{axis}
%
% relation between FT and FS
%
%
% omega between FT and FS
%
\node at (22/40*\paperwidth,6/8*\paperheight) [align=center,color=\FCOLOR] {\Huge ${\omega}$};
%
% t between FT and FS
%
\node at (18/40*\paperwidth,6/8*\paperheight) [align=center,color=\TCOLOR] {\Huge ${t}$};
%
% open arrow from FT to FS
%
% first third
%
\draw[C1,thick] (17/40*\paperwidth,25/32*\paperheight) -- (19/40*\paperwidth,25/32*\paperheight);
%
% second third (30 ° open)
%
% sqrt(3) approx 1.732050808
%
\draw[C1,thick] (19/40*\paperwidth,25/32*\paperheight) -- (19/40*\paperwidth+1.732050808/2*2/40*\paperwidth,25/32*\paperheight+1/2*2/40*\paperwidth);
%
% last third
%
\draw[C1,thick,-latex] (21/40*\paperwidth,25/32*\paperheight) -- (23/40*\paperwidth,25/32*\paperheight) node[above]{Sampling};
%
% impulse comb (open arrow)
%
\node at (22/40*\paperwidth,27/32*\paperheight) [align=center,color=\FCOLOR] {\Huge ${\cdot\Sha}$};
%
\node at (18/40*\paperwidth,27/32*\paperheight) [align=center,color=\TCOLOR] {\Huge ${\ast\Sha}$};
%
%
%
%
% relation between DTFT and DFT
%
%
% omega between DTFT and DFT
%
\node at (22/40*\paperwidth,5/32*\paperheight) [align=center,color=\FCOLOR] {\Huge ${\omega}$};
%
% t between DTFT and DFT
%
\node at (18/40*\paperwidth,5/32*\paperheight) [align=center,color=\TCOLOR] {\Huge ${t}$};
%
% open arrow from DTFT to DFT
%
% first third
%
\draw[C1,thick] (17/40*\paperwidth,7/32*\paperheight) -- (19/40*\paperwidth,7/32*\paperheight);
%
% second third (30 ° open)
%
% sqrt(3) approx 1.732050808
%
\draw[C1,thick] (19/40*\paperwidth,7/32*\paperheight) -- (19/40*\paperwidth+1.732050808/2*2/40*\paperwidth,7/32*\paperheight+1/2*2/40*\paperwidth);
%
% last third
%
\draw[C1,thick,-latex] (21/40*\paperwidth,7/32*\paperheight) -- (23/40*\paperwidth,7/32*\paperheight) node[above]{Sampling};
%
% dirac comb open arrow
%
\node at (22/40*\paperwidth,9/32*\paperheight) [align=center,color=\FCOLOR] {\Huge ${\cdot\Sha}$};
%
\node at (18/40*\paperwidth,9/32*\paperheight) [align=center,color=\TCOLOR] {\Huge ${\ast\Sha}$};
%
%
%
% relation between FT and DTFT
%
%
% omega between FT and DTFT
%
\node at (1/4*\paperwidth,18/40*\paperheight) [align=center,color=\FCOLOR,] {\Huge ${\omega}$};
%
% t between FT and FS
%
\node at (1/4*\paperwidth,22/40*\paperheight) [align=center,color=\TCOLOR,] {\Huge ${t}$};
%
% open arrow from FT to DTFT
%
% first third
% 7/32
\draw[C0,thick] (9/32*\paperwidth,23/40*\paperheight) -- (9/32*\paperwidth,21/40*\paperheight);
%
% second third (30 ° open)
%
% sqrt(3) approx 1.732050808
%
\draw[C0,thick] (9/32*\paperwidth,21/40*\paperheight) -- (9/32*\paperwidth+1/2*2/40*\paperheight,21/40*\paperheight-1.732050808/2*2/40*\paperheight);
%
% last third
%
\draw[C0,thick,-latex] (9/32*\paperwidth,19/40*\paperheight) -- (9/32*\paperwidth,17/40*\paperheight)  node[below]{Sampling};
%
% dirac comb open arrow
%
\node at (11/32*\paperwidth,18/40*\paperheight) [align=center,color=\FCOLOR,] {\Huge ${\ast\Sha}$};
%
\node at (11/32*\paperwidth,22/40*\paperheight) [align=center,color=\TCOLOR,] {\Huge ${\cdot\Sha}$};
%
% closed arrow from DTFT to FT
%
\draw[C0,thick,-latex] (7/32*\paperwidth,17/40*\paperheight) -- (7/32*\paperwidth,23/40*\paperheight) node[above]{Interpolation};
%
% convolution sinc closed arrow
%
\node at (5/32*\paperwidth,18/40*\paperheight) [align=center,color=\FCOLOR,] {\huge ${\cdot\rect}$};
%
% multiplication rect closed arrow
%
\node at (5/32*\paperwidth,22/40*\paperheight) [align=center,color=\TCOLOR,] {\huge ${\ast\sinc}$};
%
%
%
% relation between FS and DFT
%
%
% omega between FS and DFT
%
\node at (3/4*\paperwidth,18/40*\paperheight) [align=center,color=\FCOLOR,] {\Huge ${\omega}$};
%
% t between FS and DFT
%
\node at (3/4*\paperwidth,22/40*\paperheight) [align=center,color=\TCOLOR,] {\Huge ${t}$};
%
% closed arrow from FS to DFT
%
% first third
% 7/32
\draw[C0,thick] (25/32*\paperwidth,23/40*\paperheight) -- (25/32*\paperwidth,21/40*\paperheight);
%
% second third (30 ° open)
%
% sqrt(3) approx 1.732050808
%
\draw[C0,thick] (25/32*\paperwidth,21/40*\paperheight) -- (25/32*\paperwidth+1/2*2/40*\paperheight,21/40*\paperheight-1.732050808/2*2/40*\paperheight);
%
% last third
%
\draw[C0,thick,-latex] (25/32*\paperwidth,19/40*\paperheight) -- (25/32*\paperwidth,17/40*\paperheight)  node[below]{Sampling};
%
% dirac comb open arrow
%
\node at (27/32*\paperwidth,18/40*\paperheight) [align=center,color=\FCOLOR,] {\Huge ${\ast\Sha}$};
%
\node at (27/32*\paperwidth,22/40*\paperheight) [align=center,color=\TCOLOR] {\Huge ${\cdot\Sha}$};
%
% closed arrow from DFT to FS
%
\draw[C0,thick,-latex] (23/32*\paperwidth,17/40*\paperheight) -- (23/32*\paperwidth,23/40*\paperheight)  node[above]{Interpolation};
%
% convolution psinc closed arrow
%
\node at (21/32*\paperwidth,18/40*\paperheight) [align=center,color=\FCOLOR,] {\huge ${\cdot\rect}$};
%
% multiplication rect closed arrow
%
\node at (21/32*\paperwidth,22/40*\paperheight) [align=center,color=\TCOLOR,] {\huge ${\ast\psinc}$};
%
%
%
% arrow from DTFT to FS
%
\draw[-latex, ultra thick] (1/2*\paperwidth,1/2*\paperheight) -- node[above,rotate=\ANGLE,color=\TCOLOR,xshift=1/80*\paperwidth,anchor=south] {\FONTSIZE Periodisation + Interpolation}
node[below,rotate=\ANGLE,color=\FCOLOR,xshift=1/80*\paperwidth,anchor=north] {\FONTSIZE Sampling + Aperiodisation} (12.25/20*\paperwidth,12.25/20*\paperheight) coordinate (A);

\coordinate[below right=0.75cm*cos(\ANGLE) and 0.75cm*sin(\ANGLE) of A] (B);
\coordinate[above left=0.75cm*cos(\ANGLE) and 0.75cm*sin(\ANGLE) of A] (C);

\path[-latex] (B) edge [bend right=90,looseness=2] (C);

\draw[thick] (1/2*\paperwidth,1/2*\paperheight) -- node[above,rotate=\ANGLE,color=\TCOLOR,xshift=-1/80*\paperwidth,anchor=south] {\FONTSIZE Interpolation + Periodisation}
node[below,rotate=\ANGLE,color=\FCOLOR,xshift=-1/80*\paperwidth,anchor=north] {\FONTSIZE Aperiodisation + Sampling} (7.75/20*\paperwidth,7.75/20*\paperheight) coordinate (A);

\coordinate[below right=0.75cm*cos(\ANGLE) and 0.75cm*sin(\ANGLE) of A] (B);
\coordinate[above left=0.75cm*cos(\ANGLE) and 0.75cm*sin(\ANGLE) of A] (C);

\path[-latex] (B) edge [bend left=90,looseness=2] (C);

% arrow from FT to DFT
%
\draw[-latex, ultra thick] (1/2*\paperwidth,1/2*\paperheight) -- node[above,rotate=-\ANGLE,color=\TCOLOR,xshift=1/80*\paperwidth,anchor=south] {\FONTSIZE Sampling + Periodisation}
node[below,rotate=-\ANGLE,color=\FCOLOR,xshift=1/80*\paperwidth,anchor=north] {\FONTSIZE Periodisation + Sampling} (12.25/20*\paperwidth,7.75/20*\paperheight) coordinate (A);

\coordinate[below left=0.75cm*cos(\ANGLE) and 0.75cm*sin(\ANGLE) of A] (B);
\coordinate[above right=0.75cm*cos(\ANGLE) and 0.75cm*sin(\ANGLE) of A] (C);

\path[-latex] (B) edge [bend right=90,looseness=2] (C);;


\draw[thick] (1/2*\paperwidth,1/2*\paperheight) -- node[above,rotate=-\ANGLE,color=\TCOLOR,xshift=-1/80*\paperwidth] {\FONTSIZE Periodisation + Sampling}
node[below,rotate=-\ANGLE,color=\FCOLOR,xshift=-1/80*\paperwidth] {\FONTSIZE Sampling + Periodisation} (7.75/20*\paperwidth,12.25/20*\paperheight) coordinate (A);

\coordinate[below left=0.75cm*cos(\ANGLE) and 0.75cm*sin(\ANGLE) of A] (B);
\coordinate[above right=0.75cm*cos(\ANGLE) and 0.75cm*sin(\ANGLE) of A] (C);

\path[-latex] (B) edge [bend left=90,looseness=2] (C);;
%
%
%
\begin{comment}
\end{comment}
\end{tikzpicture}
\end{document}
